\noindent\textit{Sửa từ đề cũ số 14. Biên soạn: Trang.}

\begin{bt}[2,0 điểm]
Thực hiện phép tính:
\begin{listEX}[1]
\item $-\dfrac{2}{3}+25\%-\left|-5\dfrac{1}{4}\right|+(-2015)^{0}$
\item $0{,}6\cdot\left|-\dfrac{13}{8}\right|-\sqrt{\dfrac{25}{64}}:\left(\dfrac{1}{4}-\dfrac{3}{2}\right)^{2}$
\item $15\dfrac{2}{11}:\left(-\dfrac{7}{5}\right)+29\dfrac{2}{11}\cdot\sqrt{\dfrac{25}{49}}$
\item $\dfrac{15^{15}\cdot 5^{10}\cdot 10^{15}}{9^{7}\cdot 25^{20}\cdot 4^{8}}$
\end{listEX}
\loigiai{
a) $-\dfrac{14}{3}$.\quad
b) $\dfrac{23}{40}$.\quad
c) $10$.\quad
d) $\dfrac{3}{2}$.
}
\end{bt}

\begin{bt}[2,0 điểm]
Tìm $x$, biết:
\begin{listEX}[2]
\item $\dfrac{3}{4}x-\dfrac{1}{2}=\dfrac{1}{6}$
\item $\dfrac{7}{3}-\dfrac{1}{4}:(3x-1)=(-2)^{2}$
\item $\left|\dfrac{1}{5}-2x\right|-\left(1+\dfrac{1}{2}\right)^{2}=-2\dfrac{1}{10}$
\item $3^{2x-1}=27\cdot 3^{x}$
\end{listEX}
\loigiai{
a) $x=\dfrac{8}{9}$.\\
b) $x=\dfrac{17}{60}$, với điều kiện $3x-1\neq 0$.\\
c) $x=\dfrac{1}{40}$ hoặc $x=\dfrac{7}{40}$.\\
d) $x=4$.
}
\end{bt}

\begin{bt}[1,0 điểm]
Theo quy định của một cửa hàng điện thoại, chỉ tiêu của một nhân viên là bán được $20$ chiếc trong một tháng. Nhân viên hoàn thành chỉ tiêu nhận lương cơ bản $8\,000\,000$ đồng. Nếu bán vượt chỉ tiêu thì được thưởng thêm $7{,}5\%$ tiền lời của số điện thoại vượt chỉ tiêu. Trong tháng $10$, anh Hòa nhận $8\,900\,000$ đồng, gồm lương cơ bản và thưởng. Hỏi anh Hòa đã bán vượt chỉ tiêu bao nhiêu chiếc, biết mỗi chiếc bán ra cửa hàng lời $1\,500\,000$ đồng?
\loigiai{
$8$ chiếc.\\
Tiền thưởng là $900\,000$ đồng. Mỗi chiếc vượt chỉ tiêu được thưởng $1\,500\,000\cdot 7{,}5\%=112\,500$ đồng.
}
\end{bt}

\begin{bt}[1,0 điểm]
Một đội công nhân dự định hoàn thành một công trình trong $15$ ngày. Sau $6$ ngày làm việc, đội nhận thêm $4$ người. Năng suất mỗi người không đổi và phần việc còn lại được hoàn thành trong $6$ ngày. Hỏi lúc đầu đội có bao nhiêu công nhân? Đội đã hoàn thành công trình trong bao nhiêu ngày?
\loigiai{
Lúc đầu có $8$ công nhân; hoàn thành trong $12$ ngày.\\
Gọi $x$ là số công nhân lúc đầu ($x$ nguyên dương). Khối lượng công việc tính theo ngày-công thỏa $15x=6x+6(x+4)$.
}
\end{bt}

\begin{bt}[1,5 điểm]
Cho hai đường thẳng song song $a$ và $b$. Cát tuyến $c$ cắt $a$ tại $A$, cắt $b$ tại $B$. Biết $\widehat{A_{1}}=60^{\circ}$ là góc trong phía bên phải cát tuyến. Qua $A$ vẽ đường thẳng $d$ vuông góc với $a$.
\begin{listEX}[1]
\item Tính góc so le trong $\widehat{B_{1}}$ của $\widehat{A_{1}}$.
\item Vẽ tia $Bt$ là tia phân giác của $\widehat{B_{1}}$. Tính góc giữa tia $Bt$ và đường thẳng $b$.
\item Chứng minh $d\perp b$.
\end{listEX}
\begin{center}
\begin{tikzpicture}[scale=1.15,line cap=round]
	\coordinate (A) at (2,2);
	\coordinate (B) at (3.155,0);
	\coordinate (R) at (3.5,2);
	\coordinate (Up) at (2,2.85);
	\coordinate (C1) at (1.725,2.476);
	\coordinate (C2) at (3.355,-0.346);
	\draw (0,2) -- (4.8,2) node[right] {$a$};
	\draw (0,0) -- (4.8,0) node[right] {$b$};
	\draw (C1) -- (C2);
	\node[above left] at (C1) {$c$};
	\draw (Up) -- (2,1.05);
	\node[left] at (2,2.45) {$d$};
	\node[above left] at (A) {$A$};
	\node[below] at (B) {$B$};
	\pic[draw,angle radius=6mm,angle eccentricity=1.45,"$60^{\circ}$"] {angle = B--A--R};
	\pic[draw,angle radius=3.2mm] {right angle = Up--A--R};
\end{tikzpicture}
\end{center}
\loigiai{
a) $\widehat{B_{1}}=60^{\circ}$.\\
b) $30^{\circ}$.\\
c) $d\perp a$ và $a\parallel b$ nên $d\perp b$.
}
\end{bt}

\begin{bt}[2,0 điểm]
Cho đoạn thẳng $AB$. Hai điểm $M$ và $N$ không thuộc đường thẳng $AB$, $M$ khác $N$, thỏa mãn $MA=MB$ và $NA=NB$. Gọi $H$ là giao điểm của $MN$ và $AB$. Lấy điểm $C$ trên đường thẳng $MN$ sao cho $C$ khác $H$.
\begin{listEX}[1]
\item Chứng minh đường thẳng $MN$ là đường trung trực của $AB$.
\item Chứng minh $H$ là trung điểm của $AB$ và $MN\perp AB$.
\item Chứng minh $\triangle AMC=\triangle BMC$ và tia $CM$ là tia phân giác của $\widehat{ACB}$.
\end{listEX}
\loigiai{
a) $M$ và $N$ cùng cách đều hai đầu $A$, $B$ nên cả hai điểm thuộc đường trung trực của $AB$. Hai điểm phân biệt xác định đúng đường trung trực đó.\\
b) Giao điểm của đường trung trực với $AB$ là trung điểm $H$ và $MN\perp AB$.\\
c) Điểm $C$ thuộc đường trung trực nên $CA=CB$. Kết hợp $MA=MB$ và $CM$ chung, ta có $\triangle AMC=\triangle BMC$ (c.c.c). Suy ra $\widehat{ACM}=\widehat{BCM}$.
}
\end{bt}

\begin{bt}[0,5 điểm]
Tìm số tự nhiên $x>1$ thỏa mãn
\[
\dfrac{2}{2\cdot 3}+\dfrac{2}{3\cdot 4}+\dfrac{2}{4\cdot 5}+\cdots+\dfrac{2}{x(x+1)}=\dfrac{2013}{2015}.
\]
\loigiai{
$x=2014$.\\
\textit{Hướng dẫn.} $\dfrac{2}{k(k+1)}=2\left(\dfrac{1}{k}-\dfrac{1}{k+1}\right)$. Tổng từ $k=2$ đến $k=x$ bằng $\dfrac{x-1}{x+1}$.
}
\end{bt}
